% Part: second-order-logic
% Chapter: sol-and-sets
% Section: power-of-continuum

\documentclass[../../../include/open-logic-section]{subfiles}

\begin{document}

\olfileid[fr]{sol}{set}{pow}

\olsection{La puissance du continu}

\begin{explain}
La logique du second ordre permet de quantifier
sur les sous-ensembles du !!{domain}, mais pas
sur les ensembles de tels sous-ensembles. Pour
le faire directement, il faudrait une logique
\emph{du troisième ordre}. Supposons, par
exemple, que nous voulions énoncer le théorème
de Cantor : aucune fonction !!{injective} ne va
de l'ensemble des parties d'un ensemble vers
cet ensemble. Nous pourrions tenter d'écrire :
« Pour tous ensembles~$X$ et~$P$, si $P$ est
l'ensemble des parties de~$X$, alors
$\cardle{P}{X}$ est faux. » Il faudrait exprimer
que les !!{element}s de~$P$ sont exactement les
sous-ensembles de~$X$, par exemple au moyen de
$\lforall[Y][(P(Y)\liff Y\subseteq X)]$.
Mais $P(Y)$ n'est pas une !!{formula} du second
ordre : une !!{variable} de relation unaire
comme~$P$ ne prend que des termes pour arguments.

Nous pouvons toutefois \emph{simuler} une
quantification sur des ensembles d'ensembles
si le !!{domain} est assez grand. Une relation
binaire~$R$ met en rapport des !!{element}s du
!!{domain}. Pour $x\in\Domain{M}$, l'ensemble
$\Setabs{y\in\Domain{M}}{R(x,y)}$ est celui que
$x$ \emph{code par~$R$}. Réciproquement, si
$Z\subseteq\Pow{\Domain{M}}$ et que
$\Domain{M}$ contient au moins autant
d'!!{element}s que~$Z$ contient d'ensembles,
une relation $R\subseteq\Domain{M}^2$ peut
faire coder chaque $Y\in Z$ par un certain~$x$.
\end{explain}

\begin{defn}
Pour $R\subseteq\Domain{M}^2$, nous disons
que \emph{$x$ code par~$R$} l'ensemble
$\Setabs{y\in\Domain{M}}{R(x,y)}$.
\end{defn}

Si un !!{element}~$x\in\Domain{M}$ code
par~$R$ un sous-ensemble de~$\Domain{M}$,
un ensemble $Y\subseteq\Domain{M}$ code
une collection d'ensembles : ceux que codent
ses !!{element}s. Il peut ainsi coder par~$R$
l'ensemble des parties de~$X$. Cela signifie que
chaque $Z\subseteq X$ est codé par un
$x\in Y$, et que chaque $x\in Y$ code
un sous-ensemble de~$X$.

\begin{prop}
La !!{formula}
\[
\fn{Codes}(x,R,Z)\ident
\lforall[y][(Z(y)\liff R(x,y))]
\]
exprime que $s(x)$ code $s(Z)$ par~$s(R)$.
La !!{formula}
\begin{multline*}
\fn{Pow}(Y,R,X)\ident{}\\
\lforall[Z][(Z\subseteq X\lif
\lexists[x][(Y(x)\land\fn{Codes}(x,R,Z))])]\land{}\\
\lforall[x][(Y(x)\lif
\lforall[Z][(\fn{Codes}(x,R,Z)\lif Z\subseteq X)])]
\end{multline*}
exprime que $s(Y)$ code par~$s(R)$
l'ensemble des parties de~$s(X)$ : ses !!{element}s
codent exactement les sous-ensembles de~$s(X)$.
\end{prop}

\begin{explain}
Ce codage permet d'exprimer des assertions sur
les ensembles de parties en quantifiant sur
leurs codes. Le théorème de Cantor, par exemple,
dit qu'aucune fonction !!{injective} ne peut
envoyer l'ensemble des codes~$Y$ de
$\Pow{X}$ dans~$X$.
\end{explain}

\begin{prop}
L'!!{sentence}
\begin{multline*}
\lforall[X][\lforall[Y][\lforall[R][(
\fn{Pow}(Y,R,X)\lif{}\\
\lnot\lexists[u][(
\lforall[x][\lforall[y][((Y(x)\land Y(y)\land
\eq[u(x)][u(y)])\lif\eq[x][y])]]\land{}\\
\lforall[x][(Y(x)\lif X(u(x)))])])]]]
\end{multline*}
est valide.
\end{prop}

\begin{explain}
L'ensemble des parties d'un ensemble
!!{denumerable} est !!{nonenumerable} et
plus grand que tout ensemble de cardinal~$\aleph_0$.
La taille de $\Pow{\Nat}$ est appelée la
\emph{puissance du continu}, car c'est aussi
la taille de la droite réelle~$\Real$.
Si le !!{domain} est assez grand pour coder
l'ensemble des parties d'un ensemble~$X$ de
cardinal~$\aleph_0$, un ensemble~$Y$ a la
taille du continu lorsqu'il correspond
bijectivement aux sous-ensembles de~$X$ qu'il
code. Si le !!{domain} est trop petit pour
contenir un sous-ensemble équipotent à~$\Real$,
aucune relation n'y code ainsi tout
$\Pow{X}$ pour un tel~$X$.
\end{explain}

\begin{prop}
Si $\cardle{\Real}{\Domain{M}}$, la !!{formula}
\begin{multline*}
\fn{Cont}(Y)\ident
\lexists[X][\lexists[R][(
\fn{Aleph}_0(X)\land\fn{Pow}(Y,R,X)\land{}\\
\lforall[x][\lforall[y][((Y(x)\land Y(y)\land
\lforall[z][(R(x,z)\liff R(y,z))])
\lif\eq[x][y])]])]]
\end{multline*}
exprime que $\cardeq{s(Y)}{\Real}$.
\end{prop}

\begin{proof}
$\fn{Pow}(Y,R,X)$ dit que $s(Y)$ code
par~$s(R)$ tous les sous-ensembles de~$s(X)$,
et $\fn{Aleph}_0(X)$ dit que $s(X)$ est
!!{denumerable}. L'ensemble $s(Y)$ a donc
au moins la taille du continu, mais il
pourrait être plus grand si plusieurs de ses
!!{element}s codaient le même sous-ensemble
de~$s(X)$. Le dernier conjoint l'exclut :
il rend !!{injective} l'association qui envoie
un !!{element} de~$s(Y)$ sur l'ensemble
qu'il code par~$s(R)$. Réciproquement, si $s(Y)$ a la taille du continu, choisissons un sous-ensemble dénombrable infini $X_0$ du domaine et une bijection $b\colon s(Y)\to\Pow{X_0}$. Définissons $R_0$ par $R_0(a,z)$ si et seulement si $a\in s(Y)$ et $z\in b(a)$, les lignes hors de~$s(Y)$ étant vides. Sous la sémantique standard, $X_0$ et $R_0$ fournissent les témoins de $\fn{Cont}(Y)$.
\end{proof}

\begin{prop}
$\cardeq{\Domain{M}}{\Real}$ si et seulement si
\begin{multline*}
\Sat{M}{\lexists[Y][(\fn{Cont}(Y)\land
\lexists[u][(
\lforall[x][Y(u(x))]\land{}\\
\lforall[x][\lforall[y][(\eq[u(x)][u(y)]
\lif\eq[x][y])]]\land{}\\
\lforall[y][(Y(y)\lif
\lexists[x][\eq[y][u(x)]])])])]}.
\end{multline*}
\end{prop}

\begin{explain}
L'\emph{hypothèse du continu} affirme que la
taille du continu est le premier cardinal
!!{nonenumerable} : autrement dit, que
$\Pow{\Nat}$ est de cardinal~$\aleph_1$.
\end{explain}

\begin{prop}
Dans le cadre usuel avec choix, l'hypothèse du
continu est vraie si et seulement si l'!!{sentence}
\[
\fn{CH}\ident
\lforall[X][(\fn{Aleph}_1(X)\liff\fn{Cont}(X))]
\]
est valide sous la sémantique standard.
\end{prop}

La validité de $\lnot\fn{CH}$ n'équivaut pas
à la fausseté de l'hypothèse du continu.
Dans un !!{domain} !!{enumerable}, aucun
sous-ensemble n'a la taille de~$\aleph_1$
ni celle du continu; $\fn{CH}$ y est donc
toujours vrai. On peut cependant construire
un autre !!{sentence} dont la validité
équivaut à la négation de l'hypothèse.

\begin{prop}
Dans le cadre usuel avec choix, l'hypothèse du
continu est fausse si et seulement si l'!!{sentence}
\[
\fn{NCH}\ident
\lforall[X][(\fn{Cont}(X)\lif
\lexists[Y][(Y\subseteq X\land
\lnot\fn{Count}(Y)\land
\lnot\cardeq{X}{Y})])]
\]
est valide sous la sémantique standard.
\end{prop}

\begin{editorial}
La seconde formule de la première proposition
comporte dans la source une parenthèse
manquante à la fin de la clause qui porte
sur les codes; nous rétablissons sa portée.

La formule source du théorème de Cantor exige
que la fonction~$u$ soit !!{injective} sur
tout le !!{domain}, alors que le théorème
concerne sa restriction aux codes~$Y$. Nous
limitons l'injectivité à~$Y$, comme dans la
section précédente.

La formule source censée caractériser
$\cardeq{\Domain{M}}{\Real}$ demande
seulement que l'image d'une injection du
!!{domain} \emph{contienne}~$Y$ : l'identité
la satisfait dès qu'un ensemble de codes
du continu existe, même si le !!{domain}
est plus grand que le continu. Nous employons
$\fn{Cont}(Y)$, qui exclut les codes en double,
et une bijection explicite du !!{domain}
sur~$Y$. Dans la preuve de $\fn{Cont}(Y)$,
la source écrit aussi $s(Z)$ là où il faut
$s(X)$. Nous explicitons aussi le sens réciproque de cette proposition.

Les deux dernières équivalences avec
l'hypothèse du continu utilisent les
caractérisations des cardinaux sous l'axiome
du choix et la sémantique standard du second
ordre. Elles ne tranchent pas la valeur de
vérité de l'hypothèse du continu.
\end{editorial}

\end{document}
